% paper/macros.tex
% Late-binding citation macros for HalluLens paper.
%
% Every number in prose resolves at build time through these macros.
% Values are defined in paper/generated/values.tex via \csname lookup;
% build_numbers.py produces that file from paper/data/*.csv and
% paper/generated/figures/*.numbers.csv.
%
% Usage:
%   \result{csv_stem}{key}[precision]
%   \resdelta{csv_stem}{key_a}{key_b}[precision]
%   \resratio{csv_stem}{key_a}{key_b}[precision]
%   \resultCI{csv_stem}{mean_key,lo_key,hi_key}[precision]
%   \resultPM{csv_stem}{mean_key,err_key}[precision]
%
% Where:
%   csv_stem  - stem of the CSV file (e.g. "baseline" for baseline.csv,
%               "fig.transfer_llama_linear" for the sidecar of that figure)
%   key       - colon-joined coordinate per the CSV's key_schema
%   precision - number of decimal places for siunitx rounding (optional;
%               default is the CSV's default_precision, embedded as
%               hl@<stem>@__default_precision__ in values.tex)

\usepackage{siunitx}
\usepackage{xfp}
\usepackage{etoolbox}

% --------------------------------------------------------------------------
% All definitions below use @ in macro names (\@ifundefined, \hl@..., \@nil).
% Outside a .sty file, @ has catcode 12 (other), so we must temporarily make
% it a letter via \makeatletter / \makeatother for the definitions to parse.
% --------------------------------------------------------------------------
\makeatletter

% --------------------------------------------------------------------------
% Internal helper: look up \csname hl@<stem>@<key>\endcsname
% Raises \PackageError if undefined.
% --------------------------------------------------------------------------
\newcommand{\hlget}[2]{%
  \@ifundefined{hl@#1@#2}{%
    \PackageError{hallulens-macros}%
      {Undefined result key '#2' in CSV stem '#1'.
       Check paper/data/#1.csv and re-run 'make values'.}%
      {The key '#2' was not found in the values generated from '#1'.
       Valid keys are listed in paper/generated/values.tex.}%
    \textbf{??}%
  }{%
    \csname hl@#1@#2\endcsname%
  }%
}

% --------------------------------------------------------------------------
% Helper: get default precision for a CSV stem (falls back to 3)
% --------------------------------------------------------------------------
\newcommand{\hlgetdefprec}[1]{%
  \@ifundefined{hl@#1@__default_precision__}{3}{\csname hl@#1@__default_precision__\endcsname}%
}

% --------------------------------------------------------------------------
% \result{csv}{key}[precision]
% --------------------------------------------------------------------------
\DeclareRobustCommand{\result}[2]{%
  \@ifnextchar[%]
    {\hl@result@prec{#1}{#2}}%
    {\hl@result@prec{#1}{#2}[\hlgetdefprec{#1}]}%
}
\def\hl@result@prec#1#2[#3]{%
  \num[round-mode=places,round-precision=#3]{\hlget{#1}{#2}}%
}

% --------------------------------------------------------------------------
% \resdelta{csv}{key_a}{key_b}[precision]
% Computes key_a - key_b using \fpeval from xfp.
% (Named \resdelta, not \delta, to avoid clashing with the Greek letter.)
% --------------------------------------------------------------------------
\DeclareRobustCommand{\resdelta}[3]{%
  \@ifnextchar[%]
    {\hl@delta@prec{#1}{#2}{#3}}%
    {\hl@delta@prec{#1}{#2}{#3}[\hlgetdefprec{#1}]}%
}
\def\hl@delta@prec#1#2#3[#4]{%
  \edef\hl@va{\hlget{#1}{#2}}%
  \edef\hl@vb{\hlget{#1}{#3}}%
  \num[round-mode=places,round-precision=#4,explicit-sign]{\fpeval{\hl@va - \hl@vb}}%
}

% --------------------------------------------------------------------------
% \resratio{csv}{key_a}{key_b}[precision]
% Computes key_a / key_b using \fpeval.
% --------------------------------------------------------------------------
\DeclareRobustCommand{\resratio}[3]{%
  \@ifnextchar[%]
    {\hl@ratio@prec{#1}{#2}{#3}}%
    {\hl@ratio@prec{#1}{#2}{#3}[\hlgetdefprec{#1}]}%
}
\def\hl@ratio@prec#1#2#3[#4]{%
  \edef\hl@va{\hlget{#1}{#2}}%
  \edef\hl@vb{\hlget{#1}{#3}}%
  \num[round-mode=places,round-precision=#4]{\fpeval{\hl@va / \hl@vb}}%
}

% --------------------------------------------------------------------------
% \resultCI{csv}{mean_key,lo_key,hi_key}[precision]
% Renders as: mean [lo, hi]
% The three keys are comma-separated in a single brace group.
% --------------------------------------------------------------------------
\DeclareRobustCommand{\resultCI}[2]{%
  \@ifnextchar[%]
    {\hl@ci@prec{#1}{#2}}%
    {\hl@ci@prec{#1}{#2}[\hlgetdefprec{#1}]}%
}
\def\hl@ci@prec#1#2[#3]{%
  % Split #2 on commas: mean, lo, hi
  \hl@splitthree{#2}{\hl@keyA}{\hl@keyB}{\hl@keyC}%
  \num[round-mode=places,round-precision=#3]{\hlget{#1}{\hl@keyA}}%
  \ [\num[round-mode=places,round-precision=#3]{\hlget{#1}{\hl@keyB}},
     \num[round-mode=places,round-precision=#3]{\hlget{#1}{\hl@keyC}}]%
}

% --------------------------------------------------------------------------
% \resultPM{csv}{mean_key,err_key}[precision]
% Renders as: mean +/- err
% --------------------------------------------------------------------------
\DeclareRobustCommand{\resultPM}[2]{%
  \@ifnextchar[%]
    {\hl@pm@prec{#1}{#2}}%
    {\hl@pm@prec{#1}{#2}[\hlgetdefprec{#1}]}%
}
\def\hl@pm@prec#1#2[#3]{%
  \hl@splittwo{#2}{\hl@keyA}{\hl@keyB}%
  \num[round-mode=places,round-precision=#3]{\hlget{#1}{\hl@keyA}}%
  $\pm$%
  \num[round-mode=places,round-precision=#3]{\hlget{#1}{\hl@keyB}}%
}

% --------------------------------------------------------------------------
% Utility: split comma-delimited pairs and triples into named macros.
% --------------------------------------------------------------------------
\def\hl@splittwo#1#2#3{%
  \hl@splittwoaux#1\@nil{#2}{#3}%
}
\def\hl@splittwoaux#1,#2\@nil#3#4{%
  \def#3{#1}\def#4{#2}%
}

\def\hl@splitthree#1#2#3#4{%
  \hl@splitthreeaux#1\@nil{#2}{#3}{#4}%
}
\def\hl@splitthreeaux#1,#2,#3\@nil#4#5#6{%
  \def#4{#1}\def#5{#2}\def#6{#3}%
}

\makeatother
